\chapter{Conclusion}
\label{ch:conclusion}

This project built a teleoperation system for two seven-degree-of-freedom arms
mounted on a worn backpack frame, commanded from an instrumented mannequin
master, with a safety architecture appropriate to a manipulator standing
beside a person and a shared-autonomy layer above it.

\section{Against the objectives}

\emph{A master, built and characterised.} Fourteen potentiometer channels, two
inertial sensors, two force sensors and two buttons on a single
microcontroller, emitting a text frame at \SI{100}{\hertz}. Characterisation
found six of fourteen channels usable and reversed two long-standing verdicts,
both of which had come from computing statistics across transmitted rows
rather than distinct sensor updates.

\emph{A mapping that degrades predictably.} Position in spherical form with
each component taken from the sensor that observes it best. Vertical and
fore-and-aft motion track correctly on both arms. Lateral does not, and the
measurements show why no axis map or recalibration can fix it. The cost of
each channel reduction is quantified on \num{20430} recorded frames, and the
endpoint of the ladder reports position unavailable rather than inventing one.

\emph{A safety architecture validated by injection.} Fourteen of fourteen
injected faults handled. The two that took several attempts both failed for
the same underlying reason, which is that a mechanism was keyed on the wrong
observable: a dead-man on message arrival cannot detect a source that keeps
publishing stale data.

\emph{A workspace characterisation.} The two arms share no reachable
workspace under the as-built mount, which blocks inter-arm handover for
geometric reasons no software change can address. Asked four separate ways,
the working space in front of the wearer's own chest is unavailable at any
surface height or distance, and deleting the wearer's arms from the model
moves the boundary by at most \SI{75}{\milli\metre}, which identifies the
torso rather than the limbs as the constraint. The quantified mount
modification was subsequently applied as one deliberate pass, raising the
clearance ceiling from \SI{0.1610}{\metre} to \SI{0.2202}{\metre} and bringing
the task's placement targets from \SI{0.530}{\metre} to \SI{0.290}{\metre} off
the centre line.

\emph{A verification methodology.} Ten repeats over densified paths rather
than single calls at endpoints, and an instrument validated against a known
answer before a surprising measurement taken with it is reported. The
five-task verification reproduced exactly across runs, which is the first
reachability figure in this project to do so.

\emph{Motion generation and transfer.} The per-joint step limiter that
constituted the entire motion generation was replaced with a jerk-limited,
phase-synchronised generator, removing \SI{51.3}{\milli\metre} of unrequested
Cartesian excursion on a platform whose capture gate is
\SI{30}{\milli\metre}, and bringing peak commanded joint speed within the
manufacturer's declared limit for the first time. The simulation-to-hardware
gap was measured rather than assumed: every joint parks \SI{0.305}{\degree}
short on the side it approached from, which one scalar per arm removes three
quarters of.

\emph{Measurement in place of declaration.} A calibration stage measures the
working surface and the objects on it, recovering six objects with every
centre within \SI{1}{\milli\metre}, and hands the result to a planner that had
previously known about nothing except the wearer.

\emph{Identification of what blocks progress.} Seven items, listed in
\cref{ch:discussion} in the order they should be addressed. The first is two
soldered joints.

\section{What was not achieved}

No result in this thesis was validated on a physical Kinova arm carrying a
load. One session drove a single arm from its own camera and produced three
findings, all of them about calibration and none of them a performance result
(\cref{sec:hardware}). The task set is verified geometrically and has never
been performed with a wearer in the harness. The two-person study
the architecture was designed around was not run, and could not have been:
the direct teleoperation baseline depends on channels that are currently
incoherent.

The planning report's haptic objective was not delivered and, as
\cref{ch:discussion} argues, is not deliverable on this platform through the
available control interface.

\section{Further work}

Three of the findings above point at work that is cheap relative to what it
would unblock, and they are worth separating from the general list because
each of them is a measurement rather than a build.

\emph{Two soldered joints.} Left j2 and j4 are named by the observability
regression rather than chosen: with both frozen the left arm retains no
observable radial dimension at all, and repairing any other channel does not
restore one. They are the reason the direct-teleoperation baseline of
\cref{app:study} would not be representative of a working master, and
therefore the reason the study's central comparison could not be made.

\emph{One hand-eye calibration.} The transform between the wrist camera and
the arm is wrong by roughly \SI{8.5}{\degree} in a way that rotates with the
wrist, which is larger than the entire capture gate at working distance
(\cref{fig:budget}). It is the largest unmeasured term in the positioning
budget and the reason no open-loop grasp planned through it has landed inside
that gate. It needs the hardware and a printed target.

\emph{One clearance check that measures link geometry rather than link
origins.} At present the check samples link origins, so it returns a
comfortable margin for a pose in which the arm is against the mannequin
(\cref{sec:validity}). This is the only item on the list that is a safety
defect rather than an incompleteness, and it should be closed before anybody
wears the rig.

Beyond those, the substantive open questions are hardware ones. The right
arm's home joint angles have never been read from the machine, and the
workspace geometry depends on them; making the two arms' reachable sets meet
requires re-parking that arm mechanically, because \cref{ch:workspace}
establishes that no control or planning change can do it. The safety layer
has never run against a real arm, and the Kortex session-recovery path is the
part a real session meets first. And the exchange rate between operator
performance and wearer trust --- the question this architecture was built to
ask --- remains open, because it needs two people and an approved protocol
rather than more engineering.

\section{Closing}

The most transferable outcome of this work is not the mapping or the safety
architecture. It is the discipline that produced them. Twenty-six times a
result that looked like a system failure turned out to be a defect in the
measuring tool, and fifteen of those were caught only because something else
already had the answer. On a robot that stands beside a person, the
difference between a measurement and an artefact is not an academic
distinction, and the habit of establishing which one is in hand before acting
on it is the part of this project most worth carrying forward.
